% Generated from validated Article Draft Result. Do not hand-edit; revise the structured draft instead.
\noindent\textbf{6月末の研究では、agent runtimeの外側に安全kernelを置く設計、長いcontextで探索が早期終了するcontext rot、memory system比較を歪めるhidden confoundが相次いで論点になった。共通するのは、長時間agentの品質をmodel能力だけで説明できないという問題意識である。}\autocite{src-246f07f416dc,src-644589094c65,src-8603fd60d637}

The Unfireable Safety Kernelの著者らは、agent自身が変更できるruntimeの外側にexecution-time safety kernelを置く設計を提案した。Rust prototypeとformal/implementation checksで制御propertyを検討しており、alignmentをmodel内部だけでなく実行時の外部controlとして扱う。\autocite{src-644589094c65}


Diagnosing and Mitigating Context Rotの著者らは、long-horizon searchでcontextが長くなるにつれてpremature terminationが増えるfailure modeを報告した。複数model・benchmarkとcontext-management手法を比較し、単純にcontext windowを長くするだけでは探索の持続性を保証できないことを問題化している。\autocite{src-8603fd60d637}


MemDeltaの著者らはagent memoryの比較で、model、embedding、retrieval条件の違いがmemory system自体の差と混ざることを指摘した。componentを一つずつ変えるcontrolled baselineを重視し、rankingが実装以外の条件で反転し得るというevaluation confoundを示している。\autocite{src-246f07f416dc}


3本は対象が異なるが、長時間agentを成立させるために「実行時の外部control」「taskを途中で止めないcontext運用」「比較条件を固定したmemory評価」というmodel外の設計を求めている。agent研究の焦点が能力追加から、制御可能性と測定可能性へ広がっていることを示す。\autocite{src-246f07f416dc,src-644589094c65,src-8603fd60d637}


\begin{claimboundary}
いずれもprototypeまたはauthor-controlled benchmarkによる研究結果である。productionでの安全性・latency・memory qualityを保証するものではなく、異なるmodelやworkloadをまたぐ普遍的rankingとして扱わない。\autocite{src-246f07f416dc,src-644589094c65,src-8603fd60d637}
\end{claimboundary}
